% Supplementary movie captions (generated by make_movies.py); Movie S<k> = M<k>

\section*{Supplementary movies}

\paragraph{Movie S1 (M01\_baselines\_pristine\_vs\_precrack.mp4).} Baselines: intrinsic strength and a single flaw. The pristine zigzag sheet reaches 38.8 N/m and fails in one avalanche; a 20 \AA{} crack halves the strength (18.5 N/m) and fails by crack propagation at 9\% strain. Screened REBO2, athermal quasi-static uniaxial tension along x; colour: energy change per atom relative to the relaxed state (eV); red rings: atoms that have lost a bond; frames between stored quasi-static states are interpolated. All statements are model results.

\paragraph{Movie S2 (M02\_slit\_angle\_series.mp4).} Tilting a load-aligned slit array: three mechanisms, one minimum. Same porosity, five slit angles. 0$^\circ$: straight load paths (25.4 N/m). 20$^\circ$: tips of adjacent rows overlap and link en echelon (8.9 N/m, the minimum). 45$^\circ$: ligaments rotate before they break (18.5 N/m, modulus collapses). 60$^\circ$ and 90$^\circ$: the load crosses bridges between slits in bending (7.8 and 4.0 N/m). Screened REBO2, athermal quasi-static uniaxial tension along x; colour: energy change per atom relative to the relaxed state (eV); red rings: atoms that have lost a bond; frames between stored quasi-static states are interpolated. All statements are model results.

\paragraph{Movie S3 (M03\_tip\_overlap\_control.mp4).} Pre-registered control: tip overlap is necessary for the collapse. Three arrays at 20$^\circ$. With slits too short to overlap (porosity 0.10) the array keeps 20.0 N/m, 78\% of the untilted value; with overlapping tips (porosity 0.20 and 0.30) the ligaments between the tips fail first and link across the sheet. Screened REBO2, athermal quasi-static uniaxial tension along x; colour: energy change per atom relative to the relaxed state (eV); red rings: atoms that have lost a bond; frames between stored quasi-static states are interpolated. All statements are model results.

\paragraph{Movie S4 (M04\_hierarchy\_vs\_alignment.mp4).} Hierarchy is not a free lunch: alignment, level by level. All at porosity 0.20. The single-level square mesh (17.0 N/m) beats the nested mesh with veins both ways and round pores (12.8). Aligning the veins with the load (16.0) or elongating the fine pores along it (17.5) recovers the loss; doing both gives the strongest nested design (19.4), because each level then adds straight load paths. Screened REBO2, athermal quasi-static uniaxial tension along x; colour: energy change per atom relative to the relaxed state (eV); red rings: atoms that have lost a bond; frames between stored quasi-static states are interpolated. All statements are model results.

\paragraph{Movie S5 (M05\_alignment\_sweep\_single\_level.mp4).} Load-path alignment in a single-level mesh. Elongated pores rotated from across the load (alignment index A = $-$0.24, 8.6 N/m) through 30$^\circ$ and 60$^\circ$ to along the load (A = +0.35, 19.0 N/m): strength rises monotonically with the alignment of the material with the load, at constant mass. Screened REBO2, athermal quasi-static uniaxial tension along x; colour: energy change per atom relative to the relaxed state (eV); red rings: atoms that have lost a bond; frames between stored quasi-static states are interpolated. All statements are model results.

\paragraph{Movie S6 (M06\_disorder\_and\_gradients\_are\_costs.mp4).} Disorder and gradients are costs. Same porosity as the ordered fine mesh (11.7 N/m for this lattice registry, 11.7 to 13.3 across registries): positional jitter of 2 \AA{} (11.8 for this seed, 10.4 to 11.8 across seeds), pore-size dispersion (11.3), a Voronoi network (14.4 for this seed, 9.2 to 14.4 across seeds) and a pore-size gradient along the load (10.1, set by its weakest band). None is a toughening mechanism in this model. Screened REBO2, athermal quasi-static uniaxial tension along x; colour: energy change per atom relative to the relaxed state (eV); red rings: atoms that have lost a bond; frames between stored quasi-static states are interpolated. All statements are model results.

\paragraph{Movie S7 (M07\_flaw\_halo.mp4).} A halo of pores around a hole enlarges the flaw. A bare 20 \AA{} hole (19.1 N/m), the same hole surrounded by two rings of pores (16.2 N/m) and by far-field pores (13.2 N/m): the rings and the pores turn the hole into a wider effective flaw bounded by thin ligaments; they do not shield it. Screened REBO2, athermal quasi-static uniaxial tension along x; colour: energy change per atom relative to the relaxed state (eV); red rings: atoms that have lost a bond; frames between stored quasi-static states are interpolated. All statements are model results.

\paragraph{Movie S8 (M08\_fracture\_modes\_progressive\_vs\_avalanche.mp4).} Fracture mode is a load-transfer property. Eight parallel strips (armchair slit array) fail row by row and keep carrying 8 to 10 N/m to 25\% strain; three wide rows with 15\% width dispersion cascade in one avalanche at 14\% strain. Many parallel paths of moderate dispersion share the released load; few wide paths cannot. Screened REBO2, athermal quasi-static uniaxial tension along x; colour: energy change per atom relative to the relaxed state (eV); red rings: atoms that have lost a bond; frames between stored quasi-static states are interpolated. All statements are model results.

\paragraph{Movie S9 (M09\_closeup\_en\_echelon\_linking.mp4).} Close-up: en-echelon linking of overlapping slit tips (20$^\circ$). The short inclined ligaments between the tips of slits in adjacent rows carry the load in combined tension and shear and fail first; each failed ligament hands its load to the next one along the same diagonal, so the cracks link across the sheet in a staircase while the remaining strips keep carrying about 8 N/m. Screened REBO2, athermal quasi-static uniaxial tension along x; colour: energy change per atom relative to the relaxed state (eV); red rings: atoms that have lost a bond; frames between stored quasi-static states are interpolated. All statements are model results.

\paragraph{Movie S10 (M10\_closeup\_aligned\_hierarchy.mp4).} Close-up: the strongest nested design (veins and pores along the load). Fine pores elongated along the load inside veins that run along the load: 19.4 N/m at porosity 0.21. First damage appears at the bridges between the fine pores, the peak follows, and the fine-pore rows fail in a few steps while the veins carry the remaining load. A holdout that the data-driven predictor underestimated by a third. Screened REBO2, athermal quasi-static uniaxial tension along x; colour: energy change per atom relative to the relaxed state (eV); red rings: atoms that have lost a bond; frames between stored quasi-static states are interpolated. All statements are model results.
